\section{Refinement cancellation and the cost of extrapolation}
\label{f:extrapolation}
Coarsening a packet atlas combines exact components of one description. A different use of refinement compares approximations made at several scales. The distinction matters. An algebraic cancellation of leading error terms requires an expansion in a common vector space, and a useful numerical result also requires control of the remaining error and of the sample perturbations.

Let $E$ be a real normed vector space and let $f(h)\in E$ be defined for $0<|h|<\delta$. Fix a refinement ratio $b>1$. If $f(h)$ represents operators on different Hilbert spaces, comparison maps must first place them in a common space where the following additions are meaningful. The formulas alone do not construct those maps.

The denominators below are positive. Each level is a finite linear combination of the original samples at successively refined scales.

\begin{definition}[The actual refinement tower]
Set $T_0f(h)=f(h)$ and, for $m\geq0$, define
\[
 T_{m+1}f(h)=\frac{b^{m+1}T_mf(h/b)-T_mf(h)}{b^{m+1}-1}.
\]
\end{definition}

\begin{proposition}[Exact response of a monomial]
For $r\in\N$ and $v\in E$,
\[
 T_m(h^rv)=c_{m,r}h^rv,\qquad
 c_{m,r}=\prod_{k=1}^{m}\frac{b^{k-r}-1}{b^k-1}.
\]
The empty product equals one. In particular, $c_{m,0}=1$, and $c_{m,r}=0$ for $1\leq r\leq m$.
\end{proposition}
\begin{proof}
At level zero the assertion is the definition. If it holds at level $m$, the next recursion multiplies the coefficient by
\[
 \frac{b^{m+1}b^{-r}-1}{b^{m+1}-1}.
\]
This appends the next factor to the product. For $r=0$, each factor equals one. For $1\leq r\leq m$, the factor at $k=r$ vanishes.
\end{proof}

Suppose that a finite expansion is available,
\[
 f(h)=F+\sum_{r=1}^{m}h^rv_r+R(h).
\]
Linearity and the proposition give the exact identity
\[
 T_mf(h)-F=T_mR(h).
\]
Only a finite sum has been used. No convergence of an infinite formal series is needed. The size of the right-hand side remains an analytic question.

\Needspace{18\baselineskip}
\begin{theorem}[A quantified remainder and perturbation bound]
Let $p$ be a nonnegative integer and suppose
\[
 \norm{R(h)}\leq C|h|^p,\qquad
 \norm{g(h)-f(h)}\leq\varepsilon
 \quad(0<|h|<\delta),
\]
where $C,\varepsilon\geq0$. Define
\[
 A_{m,p}=\prod_{k=1}^{m}\frac{b^{k-p}+1}{b^k-1}.
\]
Then, throughout the same punctured interval,
\[
 \norm{T_mg(h)-F}\leq A_{m,p}C|h|^p+A_{m,0}\varepsilon.
\]
\end{theorem}
\begin{proof}
The interval is stable under $h\mapsto h/b$. The triangle inequality applied to one tower step gives
\[
 \norm{T_{m+1}R(h)}\leq
 \frac{b^{m+1}\norm{T_mR(h/b)}+\norm{T_mR(h)}}{b^{m+1}-1}.
\]
Starting with the stated remainder bound and proceeding by induction multiplies its constant at step $m+1$ by $(b^{m+1-p}+1)/(b^{m+1}-1)$. Hence $\norm{T_mR(h)}\leq A_{m,p}C|h|^p$. Apply the same argument with $p=0$ to $g-f$. Linearity and the triangle inequality separate the two contributions.
\end{proof}

When $p=m+1$, this proves the expected order improvement under an explicit remainder hypothesis. It also retains the effect of noisy samples. Each factor in $A_{m,0}$ exceeds one, so extrapolation can amplify a uniform sample error even while cancelling several deterministic powers of $h$.

\subsection{A complete three-level example}
For $b=2$, direct substitution gives
\begin{align*}
 T_1f(h)&=2f(h/2)-f(h),\\
 T_2f(h)&=\frac{8f(h/4)-6f(h/2)+f(h)}{3},\\
 T_3f(h)&=\frac{64f(h/8)-56f(h/4)+14f(h/2)-f(h)}{21}.
\end{align*}
Take the exact scalar polynomial $f(h)=1+2h-3h^2+4h^3$. The responses are
\[
 T_1f(h)=1+\frac32h^2-3h^3,\qquad
 T_2f(h)=1+\frac12h^3,\qquad T_3f(h)=1.
\]
These are polynomial identities for every real $h$. If each of the three samples used by $T_2$ has absolute error at most $\varepsilon$, its output error is at most
\[
 \frac{8+6+1}{3}\varepsilon=5\varepsilon.
\]
This bound is attained by errors with signs matching the three weights. Thus exact removal of the first two powers does not imply a contraction on arbitrary perturbations. For the displayed polynomial, the perturbed second-level approximation has the explicit total bound $|h|^3/2+5\varepsilon$.

The calculation supplies a practical reading rule for later refinement claims. An exact coefficient identity, an asymptotic remainder estimate, and a stability bound are separate pieces of evidence. A physical approximation uses all three, together with the comparison maps that make its values commensurable.


